% ============================================================
%  SECTION 4: How the exchange rate is determined (Member 2)
%  Lines marked [added at merge] are new; everything else is Member 2's text,
%  with notation made uniform (₹82 per US dollar) and figure/section numbers.
% ============================================================
\section{How the exchange rate is determined: demand and supply}\label{sec:ds}

\subsection{Finding the equilibrium exchange rate}
In \secref{sec:forex}, we saw that people and businesses demand US dollars to make payments abroad, while dollars enter the Indian foreign exchange market through sources such as export earnings and remittances. How do these demands and supplies determine the exchange rate?

We can answer this question using a demand-and-supply model. The demand curve shows the quantity of dollars that buyers wish to purchase at different exchange rates. The supply curve shows the quantity of dollars that sellers wish to supply at those rates.

% [added at merge] why the curves slope the way they do
The demand curve slopes downward. When a dollar costs fewer rupees, foreign goods, foreign education and travel abroad become cheaper for Indians, so they wish to buy more dollars. The supply curve slopes upward. When each dollar fetches more rupees, selling abroad becomes more rewarding for exporters and sending money home becomes more attractive for workers abroad, so more dollars are offered.

The two curves meet at a point called the \keyterm{equilibrium exchange rate}{The exchange rate at which the quantity of dollars demanded equals the quantity supplied.}. At this exchange rate, the quantity of dollars demanded equals the quantity supplied. There is no excess demand or excess supply at the intersection, so the model identifies no pressure for the exchange rate to change from that point.

% [added at merge] how the market gets back to equilibrium
Away from the equilibrium, there is such pressure. If the rate is above it, sellers offer more dollars than buyers want. This \keyterm{excess supply}{When more dollars are offered than buyers want at the current exchange rate; the rate tends to fall.} pushes the rate down. If the rate is below it, buyers want more dollars than are on offer. This \keyterm{excess demand}{When buyers want more dollars than are on offer at the current exchange rate; the rate tends to rise.} pushes the rate up. Either way, the rate moves back towards the equilibrium.

To represent this in a diagram, put the exchange rate, measured as rupees per US dollar (₹/\$), on the vertical axis and the quantity of dollars traded on the horizontal axis. The downward-sloping demand curve and upward-sloping supply curve intersect at the equilibrium exchange rate (Figure~\ref{fig:ds}a). The position of this intersection depends on both curves, not on either curve alone.

\begin{figure}[!htb]
  \centering
  \resizebox{\linewidth}{!}{% Figure (Section 4, Member 2): equilibrium, a rise in demand, a rise in supply.
% Redrawn from Member 2's Figure 3.1 in the shared style; values match the text
% (equilibrium ₹82, demand rise to ₹86, supply rise to ₹79).
% Geometry: D: y = 3.2 - 0.8125(x - 0.4), S: y = 0.6 + 0.8125(x - 0.4), E at (2, 1.9).
\begin{tikzpicture}[ds, x=0.75cm]
\foreach \ox/\title in {0/{(a) Initial equilibrium}, 5.7/{(b) Demand increases}, 11.4/{(c) Supply increases}} {
  \begin{scope}[shift={(\ox,0)}]
    \draw[ds axis] (0,0) -- (0,3.7) node[above] {₹ per US dollar};
    \draw[ds axis] (0,0) -- (5,0);
    \node[below] at (2.5,0) {Quantity of US dollars};
    \node[ds title] at (-0.4,4.15) {\title};
  \end{scope}
}
% (a)
\draw[ds demand] (0.4,3.2) -- (3.6,0.6) node[below] {D};
\draw[ds supply] (0.4,0.6) -- (3.6,3.2) node[above] {S};
\draw[ds guide] (0,1.9) -- (2,1.9) -- (2,0);
\node[ds eq] at (2,1.9) {}; \node[above=3pt] at (2,1.9) {\textbf{E}};
\node[left] at (0,1.9) {82};
% (b) demand shifts right by 1: E1 at (2.5, 2.306)
\begin{scope}[shift={(5.7,0)}]
  \draw[ds demand] (0.4,3.2) -- (3.6,0.6) node[below] {D\textsubscript{0}};
  \draw[ds demand, ds new] (1.4,3.2) -- (4.6,0.6) node[below] {D\textsubscript{1}};
  \draw[ds supply] (0.4,0.6) -- (3.6,3.2) node[above] {S};
  \draw[ds shift] (2.95,1.2) -- (3.75,1.2);
  \draw[ds guide] (0,1.9) -- (2,1.9) -- (2,0);
  \draw[ds guide] (0,2.306) -- (2.5,2.306) -- (2.5,0);
  \node[ds eq] at (2,1.9) {}; \node[below=3pt] at (2,1.9) {E\textsubscript{0}};
  \node[ds eq=fxgreen] at (2.5,2.306) {}; \node[above=3pt, text=fxgreen] at (2.5,2.306) {\textbf{E\textsubscript{1}}};
  \node[left] at (0,1.9) {82};
  \node[left, text=fxgreen] at (0,2.306) {\textbf{86}};
\end{scope}
% (c) supply shifts right by 0.8: E1 at (2.4, 1.575)
\begin{scope}[shift={(11.4,0)}]
  \draw[ds demand] (0.4,3.2) -- (3.6,0.6) node[below] {D};
  \draw[ds supply] (0.4,0.6) -- (3.6,3.2) node[above] {S\textsubscript{0}};
  \draw[ds supply, ds new] (1.2,0.6) -- (4.4,3.2) node[above] {S\textsubscript{1}};
  \draw[ds shift] (2.95,2.6) -- (3.55,2.6);
  \draw[ds guide] (0,1.9) -- (2,1.9) -- (2,0);
  \draw[ds guide] (0,1.575) -- (2.4,1.575) -- (2.4,0);
  \node[ds eq] at (2,1.9) {}; \node[above=3pt] at (2,1.9) {E\textsubscript{0}};
  \node[ds eq=fxgreen] at (2.4,1.575) {}; \node[below=3pt, text=fxgreen] at (2.4,1.575) {\textbf{E\textsubscript{1}}};
  \node[left] at (0,1.9) {82};
  \node[left, text=fxgreen] at (0,1.575) {\textbf{79}};
\end{scope}
\end{tikzpicture}
}
  \caption{How changes in demand and supply affect the exchange rate.}
  \label{fig:ds}
  {\raggedright\small\sffamily (a) Demand and supply meet at E, ₹82 per US dollar. (b) Demand shifts right from D\textsubscript{0} to D\textsubscript{1}; the new equilibrium E\textsubscript{1} is at ₹86, so the rupee has depreciated. (c) Supply shifts right from S\textsubscript{0} to S\textsubscript{1}; the new equilibrium is at ₹79, so the rupee has appreciated.\par}
  \figsource{Drawn by the authors. All values are hypothetical.}
\end{figure}

\subsection{What happens when demand or supply changes?}
The equilibrium exchange rate need not remain at the same level. When the demand for dollars or their supply changes, the curves shift, and the market moves towards a new equilibrium.

\minihead{An increase in the demand for dollars.}
Suppose more Indians begin travelling to the United States or purchasing foreign goods and services. They need more dollars to make these payments. At each exchange rate, buyers now wish to purchase more dollars than before. The demand curve therefore shifts to the right.

Consider a hypothetical market in which the initial equilibrium exchange rate is ₹82 per US dollar. After the increase in demand, suppose the new equilibrium is ₹86. Buyers must now pay ₹86 rather than ₹82 for each dollar. The price of the dollar in rupees has risen, which means that the rupee has depreciated against the dollar.

This is the result shown in panel (b) of Figure~\ref{fig:ds}. The supply curve remains unchanged, while the demand curve shifts to the right. The new intersection lies at a higher exchange rate.

\minihead{An increase in the supply of dollars.}
Now consider the opposite situation. Suppose Indian exporters receive more dollar payments from customers abroad, increasing the supply of dollars in the Indian foreign exchange market. At each exchange rate, sellers are willing to supply more dollars than before, so the supply curve shifts to the right.

Suppose the initial equilibrium is ₹82 per US dollar and the increased supply brings the exchange rate down to ₹79. Each dollar can now be exchanged for fewer rupees. The rupee has appreciated against the dollar. Panel (c) of Figure~\ref{fig:ds} illustrates this change: the demand curve remains unchanged, while the supply curve shifts to the right and meets it at a lower exchange rate.

\subsection{Using the model to understand exchange-rate movements}
These examples isolate one change at a time. In the first, demand increases while supply is held constant; in the second, supply increases while demand is held constant. Economists describe this method as assuming that other things remain constant. It allows us to examine the effect of one change without confusing it with other changes occurring at the same time.

% [added at merge] shift versus movement along a curve
Notice the difference between a movement along a curve and a shift of the curve. A change in the exchange rate itself moves buyers and sellers along their curves. A change in anything else, such as travel plans or export orders, shifts a whole curve.

Actual foreign exchange markets are more complicated. Demand and supply may shift together. For example, stronger demand for imports may increase the demand for dollars, while higher export earnings increase their supply. The resulting exchange rate depends on the combined effects of these changes. As the demand-and-supply model shows, we must consider both curves to understand where the new equilibrium will lie \seesec{sec:moves}.

% [added at merge] practice with the model
\begin{exercise}{Draw the shift}
For each event, draw the market for US dollars in India. Decide which curve shifts, in which direction, and what happens to the rupee. Assume nothing else changes.
\begin{enumerate}[label=(\alph*), leftmargin=2em]
  \item A boom in India's software exports to the United States.
  \item Families buy more imported gold during the festival and wedding season.
  \item A large foreign investment fund sells its Indian shares and takes the money home.
\end{enumerate}
\answers{(a) Supply of dollars shifts right; the ₹/\$ rate falls, so the rupee appreciates. (b) Demand for dollars shifts right; the rate rises, so the rupee depreciates. (c) Supply of dollars shifts left; the rate rises, so the rupee depreciates.}
\end{exercise}
